%% LyX 2.4.4 created this file.  For more info, see https://www.lyx.org/.
%% Do not edit unless you really know what you are doing.
\documentclass[english]{article}
\usepackage[T1]{fontenc}
\usepackage[latin9]{inputenc}
\usepackage{amsmath}
\usepackage{amsthm}
\usepackage{geometry}
\geometry{verbose,tmargin=1in,bmargin=1in,lmargin=1in,rmargin=1in}

\makeatletter

%%%%%%%%%%%%%%%%%%%%%%%%%%%%%% LyX specific LaTeX commands.
%% Because html converters don't know tabularnewline
\providecommand{\tabularnewline}{\\}

%%%%%%%%%%%%%%%%%%%%%%%%%%%%%% Textclass specific LaTeX commands.
\numberwithin{equation}{section}
\numberwithin{figure}{section}
\newlength{\lyxlabelwidth}      % auxiliary length 
\theoremstyle{plain}
\newtheorem*{thm*}{\protect\theoremname}

%%%%%%%%%%%%%%%%%%%%%%%%%%%%%% User specified LaTeX commands.
\usepackage{fancyhdr}
\usepackage{lastpage}
\pagestyle{fancy}
\usepackage{enumitem}
\usepackage{ifthen}
\everymath=\expandafter{\the\everymath\displaystyle}
%\DeclareMathSizes{10}{10}{10}{10}
\usepackage{multicol}

\usepackage{tikz}
\usetikzlibrary{patterns}
\usetikzlibrary{plotmarks}
\usepackage{pgfplots}
\pgfplotsset{compat=newest}

\usepackage{calc}
\usepackage[nomessages]{fp}

\usepackage{datetime}
% Added by lyx2lyx
\setlength{\parskip}{\smallskipamount}
\setlength{\parindent}{0pt}

\makeatother

\usepackage{babel}
\providecommand{\theoremname}{Theorem}

\begin{document}
\renewcommand{\author}{Dr.\,James Ashe}

\newcommand{\classNum}{335}

\newcommand{\classSec}{A}

\newcommand{\nameType}{}

\newcommand{\examType}{The dihedral group of degree 4}

\newcommand{\Date}{\formatdate{24}{5}{2013}} %%\AdvanceDate[12] \today

%\newdate{my_date}{\day}{\month}{\year}%   
%\AdvanceDate[-5] 
%\displaydate{my_date}

%%First page's header%%%%%%%%%%%%%%%%%%%%%%
\fancypagestyle{firststyle} %%header settings for first page
{
	\fancyhead[L]{MTH \classNum{}(\classSec{}) \author{}
	}
	\fancyhead[R]{\textbf{\examType{}} \hspace{1cm} \Date{}}
	\setlength{\headsep}{14pt}
}
%%%%%%%%%%%%%%%%%%%%%%%%%%%%%%%%%%%%%%%%%%

%%Next page's header%%%%%%%%%%%%%%%%%%%%%%%
%\setlist{nolistsep} %eliminates space between list items
\lhead{MTH \classNum{}(\classSec{})} 
\chead{\textbf{\nameType}} 
\cfoot{\thepage{}~of~\pageref{LastPage}} 
\setlength{\headsep}{5pt} 
%%%%%%%%%%%%%%%%%%%%%%%%%%%%%%%%%%%%%%%%%%%

%%Various settings; see the preamble for other settings%%%%%
%%\setenumerate[0]{leftmargin=12pt,itemindent=0pt,labelwidth=0pt}
\renewcommand{\labelenumi}{\textbf{\arabic{enumi}.}}
\renewcommand{\labelenumii}{\textbf{\alph{enumii}.}}
\thispagestyle{firststyle}
%%%%%%%%%%%%%%%%%%%%%%%%%%%%%%%%%%%%%%%%%%

Consider a square placed on Cartesian plane such that the square's
center is at the origin and its corners are on the axes (like a diamond;
see the figure below). Any motion of the square which keeps the center
at the origin and corners on the axes has the same result as one of
eight motions: $r_{0}\equiv\mbox{ rotation of }0^{\circ}$, $r_{1}\equiv\mbox{ rotation of }90^{\circ}$,
$r_{2}\equiv\mbox{ rotation of }180^{\circ}$, $r_{3}\equiv\mbox{ rotation of }270$,
$s_{0}\equiv\mbox{ }x-\mbox{axis reflection}$, $s_{1}\equiv\mbox{reflection on }y=x$,
$s_{2}\equiv\,y-\mbox{axis reflection}$, or , and $s_{3}\equiv\mbox{reflection on }y=-x$.

There are exactly 8 unique ways to position the square. There are
four places to position corner $1$. Corner $2$ must be adjacent
to $1$; so there are two choices for its place giving 8 possibilities.
Corners $3$ and $4$'s placement is then determined as they must
be adjacent to $2$ and $1$ respectively.
\begin{thm*}
$D_{4}=\left\{ r_{0},r_{1},r_{2},r_{3},s_{0},s_{1},s_{2},s_{3}\right\} $
is a group under composition of these motions. The set $D_{4}$ is
called the \emph{dihedral group of degree 4} (the group of symmetries
of a regular 4 sided polygon).First 
\end{thm*}
\begin{proof}
The identity is $r_{0}$. (The inverse of a rotation, $r_{a}$, is
the rotation, $r_{b}$, such that $a+b=0\bmod4$, i.e., adds up to
$2\pi$. Also $s_{n}^{-1}=s_{n}$) For associativity, note that repositioning
the corners is a bijective function from $f:\left\{ 1,2,3,4\right\} \rightarrow\left\{ 1,2,3,4\right\} $,
i.e., it is a permutation on $\left\{ 1,2,3,4\right\} $. So each
motion is a function, and function composition is associative. For
closure and inverses, see the Cayley table below:

\begin{multicols}{2}

\hspace{2cm}%
\begin{tabular}{c|ccccccccc}
 & $r_{0}$ & $r_{1}$ & $r_{2}$ & $r_{3}$ & $s_{0}$ & $s_{1}$ & $s_{2}$ & $s_{3}$ & $s_{4}$\tabularnewline
\hline 
$r_{0}$ & $r_{0}$ & $r_{1}$ & $r_{2}$ & $r_{3}$ & $s_{0}$ & $s_{1}$ & $s_{2}$ & $s_{3}$ & $s_{4}$\tabularnewline
$r_{1}$ & $r_{1}$ &  &  &  &  &  &  &  & \tabularnewline
$r_{2}$ & $r_{2}$ &  &  &  &  &  &  &  & \tabularnewline
$r_{3}$ & $r_{3}$ &  &  &  &  &  &  &  & \tabularnewline
$s_{0}$ & $s_{0}$ &  &  &  &  &  &  &  & \tabularnewline
$s_{1}$ & $s_{1}$ &  &  &  &  &  &  &  & \tabularnewline
$s_{2}$ & $s_{2}$ &  &  &  &  &  &  &  & \tabularnewline
$s_{3}$ & $s_{3}$ &  &  &  &  &  &  &  & \tabularnewline
\end{tabular}\columnbreak\hspace*{2cm}\newcommand{\xmax}{1.5}
\newcommand{\xmin}{-1.5}
\newcommand{\xstep}{0} %must be xmin+n
\newcommand{\ymax}{1.5}
\newcommand{\ymin}{-1.5}
\newcommand{\ystep}{0} %must be ymin+n
%%%%%%%
\begin{tikzpicture} 
\begin{axis}[height=3.5cm, 
	width=3.5cm, 
	axis lines=middle,
	x axis line style={-},y axis line style={-}, 
	grid=none,
	%axis line style={-latex}, 
	xmin=\xmin,xmax=\xmax, 
	ymin=\ymin,ymax=\ymax,
	%restrict y to domain=-1:10,
	no markers, 
	xlabel={$$},ylabel={$$},
	ticks=none,
	]
	
	\addplot[very thick,blue] coordinates{(1,0) (0,1)};
	\addplot[very thick,blue] coordinates{(0,1) (-1,0)};
	\addplot[very thick,blue] coordinates{(-1,0) (0,-1)};
	\addplot[very thick,blue] coordinates{(0,-1) (1,0)};

	\node at (axis cs:1.3,.3) {$1$};
	\node at (axis cs:-.3,1.3) {$2$};
	\node at (axis cs:-1.3,-.3) {$3$};
	\node at (axis cs:.3,-1.3) {$4$};
	
\end{axis}
\end{tikzpicture}\end{multicols}
\end{proof}
Alternatively,
\begin{proof}
The motions are linear transformations. So each motion can be represented
with a $2\times2$ matrix. 
\[
r_{0}=\left(\begin{array}{cc}
1 & 0\\
0 & 1
\end{array}\right),\,r_{1}=\left(\begin{array}{cc}
0 & -1\\
1 & 0
\end{array}\right),\,r_{2}=\left(\begin{array}{cc}
-1 & 0\\
0 & -1
\end{array}\right),\,r_{3}=\left(\begin{array}{cc}
0 & 1\\
-1 & 0
\end{array}\right)
\]
\[
s_{0}=\left(\begin{array}{cc}
1 & 0\\
0 & -1
\end{array}\right),\,s_{1}=\left(\begin{array}{cc}
0 & 1\\
1 & 0
\end{array}\right),\,s_{2}=\left(\begin{array}{cc}
-1 & 0\\
0 & 1
\end{array}\right),\,s_{3}=\left(\begin{array}{cc}
0 & -1\\
-1 & 0
\end{array}\right)
\]

Where a corner's Cartesian coordinates $\left(x,y\right)$ are moved
via matrix multiplication:
\[
\left(\begin{array}{cc}
a & b\\
c & d
\end{array}\right)\left(\begin{array}{c}
x\\
y
\end{array}\right)=\left(\begin{array}{c}
ax+by\\
cx+dy
\end{array}\right)
\]

\[
r_{0}\left(\begin{array}{c}
x\\
y
\end{array}\right)=\left(\begin{array}{c}
x\\
y
\end{array}\right),\,r_{1}\left(\begin{array}{c}
x\\
y
\end{array}\right)=\left(\begin{array}{c}
-y\\
x
\end{array}\right),\,r_{2}\left(\begin{array}{c}
x\\
y
\end{array}\right)=\left(\begin{array}{c}
-x\\
-y
\end{array}\right),\,r_{3}\left(\begin{array}{c}
x\\
y
\end{array}\right)=\left(\begin{array}{c}
y\\
-x
\end{array}\right)
\]
\[
s_{0}\left(\begin{array}{c}
x\\
y
\end{array}\right)=\left(\begin{array}{c}
x\\
-y
\end{array}\right),\,s_{1}\left(\begin{array}{c}
x\\
y
\end{array}\right)=\left(\begin{array}{c}
y\\
x
\end{array}\right),\,s_{2}\left(\begin{array}{c}
x\\
y
\end{array}\right)=\left(\begin{array}{c}
-x\\
y
\end{array}\right),\,s_{3}\left(\begin{array}{c}
x\\
y
\end{array}\right)=\left(\begin{array}{c}
-y\\
-x
\end{array}\right)
\]

$r_{0}$ is the identity. Matrices are associative. Each matrix is
square, and has a non-zero determinant. So it is easy to check with
$\frac{1}{ad-bc}$$\left(\begin{array}{cc}
d & -b\\
-c & a
\end{array}\right)=\left(\begin{array}{cc}
a & b\\
c & d
\end{array}\right)^{-1}$ that each element has an inverse in $D_{4}$. For closure, recall
that the $\det\left(A\right)\det\left(B\right)=\det\left(AB\right)$.
The determinant for each element in $D_{4}$ is $\pm1$. So the determinant
of the product of any two elements in $D_{4}$ will also be $\pm1$.
Since the only entries in the matrices are $0$ and $1$, $\det\left(\begin{array}{cc}
a & b\\
c & d
\end{array}\right)=ad-bc$ implies that $a=d=0$ or $b=c=0$ with the remaining values being
$\pm1$. As all 8 such possibilities are in $D_{4}$ this implies
that the product must also be in $D_{4}$.
\end{proof}

\end{document}
